% =========================================================
% 09_column_null_spaces.tex
% Vectors and Linear Algebra
% =========================================================

% =========================================================
% SECTION DIVIDER
% =========================================================

\section{Column Space and Null Space}

% =========================================================
% COLUMN SPACE
% =========================================================

\begin{frame}{Column Space}

Let

\[
A=
\begin{bmatrix}
|&|& &|\\
\vect{a}_1&\vect{a}_2&\cdots&\vect{a}_n\\
|&|& &|
\end{bmatrix}.
\]

The \textbf{column space} of $A$ is the set of all
linear combinations of its columns.

\[
\boxed{
\operatorname{Col}(A)
=
\operatorname{span}
\{\vect{a}_1,\vect{a}_2,\ldots,\vect{a}_n\}
}
\]

\vspace{0.4cm}

Therefore, the column space contains every vector that can be
obtained in the form

\[
A\vect{x}.
\]

\end{frame}


% =========================================================
% COLUMN SPACE AND Ax=b
% =========================================================

\begin{frame}{Column Space and $A\vect{x}=\vect{b}$}

Consider

\[
A\vect{x}=\vect{b}.
\]

Since

\[
A\vect{x}
=
x_1\vect{a}_1+
x_2\vect{a}_2+\cdots+x_n\vect{a}_n,
\]

a solution exists precisely when $\vect{b}$ can be expressed
as a linear combination of the columns of $A$.

Therefore,

\[
\boxed{
A\vect{x}=\vect{b}
\text{ is consistent}
\iff
\vect{b}\in\operatorname{Col}(A).
}
\]

The column space therefore describes all possible outputs of
the transformation represented by $A$.

\end{frame}


% =========================================================
% COLUMN SPACE EXAMPLE
% =========================================================

\begin{frame}{Example: Column Space}

Consider
\vspace{-0.1cm}
\[
A=
\begin{bmatrix}
1&2\\
2&4
\end{bmatrix}.
\]

Its columns are
\vspace{-0.1cm}
\[
\vect{a}_1=
\begin{bmatrix}
1\\2
\end{bmatrix},
\qquad
\vect{a}_2=
\begin{bmatrix}
2\\4
\end{bmatrix}.
\]

Since

\[
\vect{a}_2=2\vect{a}_1,
\]

the two columns do not generate the entire plane.

Thus,
\vspace{-0.1cm}
\[
\boxed{
\operatorname{Col}(A)
=
\operatorname{span}
\left\{
\begin{bmatrix}
1\\2
\end{bmatrix}
\right\}.
}
\]

The column space is a line through the origin.

\end{frame}


% =========================================================
% NULL SPACE
% =========================================================

\begin{frame}{Null Space}

The \textbf{null space} of a matrix $A$ is the set of all
vectors $\vect{x}$ satisfying

\[
A\vect{x}=\mathbf{0}.
\]

It is denoted by

\[
\boxed{
\operatorname{Null}(A)
=
\{\vect{x}:A\vect{x}=\mathbf{0}\}.
}
\]

\vspace{0.4cm}

The null space contains all inputs that are mapped to
the zero vector.

\[
\vect{x}
\xrightarrow{\quad A\quad}
\mathbf{0}.
\]

\end{frame}


% =========================================================
% NULL SPACE EXAMPLE
% =========================================================

\begin{frame}{Example: Null Space}

Consider
\vspace{-1.5cm}
\begin{columns}[c,totalwidth=\textwidth]

% ---------- LEFT ----------
\begin{column}{0.42\textwidth}

\begin{center}

\[
A=
\begin{bmatrix}
1&2\\
2&4
\end{bmatrix}
\]

\end{center}

\end{column}

% ---------- RIGHT ----------
\begin{column}{0.50\textwidth}

\begin{center}

\[
A
\begin{bmatrix}
x\\y
\end{bmatrix}
=
\begin{bmatrix}
0\\0
\end{bmatrix}
\]

\end{center}

\end{column}

\end{columns}

\vspace{0.2cm}

Solving the system gives
\vspace{-0.1cm}
\[
x+2y=0
\qquad\Longrightarrow\qquad
x=-2y.
\]

Let $y=t$. Then

\[
\begin{bmatrix}
x\\y
\end{bmatrix}
=
t
\begin{bmatrix}
-2\\1
\end{bmatrix}.
\]

Therefore,

\[
\boxed{
\operatorname{Null}(A)
=
\operatorname{span}
\left\{
\begin{bmatrix}
-2\\1
\end{bmatrix}
\right\}
}
\]

\end{frame}


% =========================================================
% NULL SPACE AS KERNEL
% =========================================================

\begin{frame}{Null Space as the Kernel}

For a linear transformation

\[
T(\vect{x})=A\vect{x},
\]

the null space is also called the \textbf{kernel}.

\[
\boxed{
\operatorname{Null}(A)
=
\ker(T)
}
\]

\vspace{0.4cm}

It contains all vectors that are mapped to zero:

\[
\boxed{
\ker(T)
=
\{\vect{x}:T(\vect{x})=\mathbf{0}\}.
}
\]

\vspace{0.4cm}

The kernel measures which directions are lost or collapsed
by the transformation.

\end{frame}


% =========================================================
% BASIS OF COLUMN SPACE
% =========================================================

\begin{frame}{Basis for the Column Space}

A basis for the column space consists of linearly independent
columns that span the column space.

Consider
\vspace{-0.1cm}
\[
A=
\begin{bmatrix}
1&2&3\\
2&4&6
\end{bmatrix}.
\]

Since
\vspace{-0.1cm}
\[
\vect{a}_2=2\vect{a}_1,
\qquad
\vect{a}_3=3\vect{a}_1,
\]

only the first column is needed.

Therefore,
\vspace{-0.1cm}
\[
\boxed{
\left\{
\begin{bmatrix}
1\\2
\end{bmatrix}
\right\}
}
\]

is a basis for $\operatorname{Col}(A)$.

\end{frame}


% =========================================================
% PIVOT COLUMNS
% =========================================================

\begin{frame}{Pivot Columns and Column Space}

Gaussian elimination can be used to identify independent
columns.

Consider

\[
A=
\begin{bmatrix}
1&2&3\\
2&4&6
\end{bmatrix}.
\]

Row reduction gives

\[
\begin{bmatrix}
1&2&3\\
0&0&0
\end{bmatrix}.
\]

There is one pivot.

\vspace{0.3cm}

The corresponding column of the \textbf{original matrix}
forms a basis for the column space.

\[
\boxed{
\operatorname{rank}(A)
=
\text{number of pivot columns}.
}
\]

\end{frame}


% =========================================================
% RANK
% =========================================================

\begin{frame}{Rank of a Matrix}

The \textbf{rank} of a matrix is the dimension of its
column space.

\[
\boxed{
\operatorname{rank}(A)
=
\dim\bigl(\operatorname{Col}(A)\bigr).
}
\]

Equivalently, the rank is the number of pivot columns
obtained during row reduction.

\vspace{0.4cm}

For an $m\times n$ matrix,

\[
\boxed{
\operatorname{rank}(A)\leq\min(m,n).
}
\]

Rank measures the number of independent directions produced
by the transformation.

\end{frame}


% =========================================================
% ROW SPACE
% =========================================================

\begin{frame}{Row Space}

The \textbf{row space} of a matrix is the span of its rows.

For

\[
A=
\begin{bmatrix}
1&2&3\\
2&4&6
\end{bmatrix},
\]

the rows are

\[
\begin{bmatrix}
1&2&3
\end{bmatrix},
\qquad
\begin{bmatrix}
2&4&6
\end{bmatrix}.
\]

Since the second row is twice the first,

\[
\boxed{
\operatorname{Row}(A)
=
\operatorname{span}
\left\{
\begin{bmatrix}
1&2&3
\end{bmatrix}
\right\}.
}
\]

The dimension of the row space is also the rank of $A$.

\end{frame}


% =========================================================
% NULLITY
% =========================================================

\begin{frame}{Nullity}

The \textbf{nullity} of a matrix is the dimension of its
null space.

\[
\boxed{
\operatorname{nullity}(A)
=
\dim\bigl(\operatorname{Null}(A)\bigr).
}
\]

For an $m\times n$ matrix,

\[
\boxed{
\operatorname{rank}(A)
+
\operatorname{nullity}(A)
=
n.
}
\]

This is the \textbf{Rank--Nullity Theorem}.

\vspace{0.4cm}

The value $n$ is the number of columns of $A$.

\end{frame}


% =========================================================
% RANK-NULLITY EXAMPLE
% =========================================================

\begin{frame}{Example: Rank--Nullity}

Consider
\vspace{-0.1cm}
\[
A=
\begin{bmatrix}
1&2&3\\
2&4&6
\end{bmatrix}.
\]

There are
\vspace{-0.1cm}
\[
n=3
\]

columns.

From row reduction,
\vspace{-0.1cm}
\[
\operatorname{rank}(A)=1.
\]

Therefore,
\vspace{-0.1cm}
\[
1+\operatorname{nullity}(A)=3.
\]

Hence,
\vspace{-0.1cm}
\[
\boxed{
\operatorname{nullity}(A)=2.
}
\]

So the matrix has one independent output direction and
two independent null directions.

\end{frame}


% =========================================================
% LINEAR DEPENDENCE
% =========================================================

\begin{frame}{Linear Dependence and the Null Space}

Suppose the columns of $A$ are
\vspace{-0.1cm}
\[
\vect{a}_1,\vect{a}_2,\ldots,\vect{a}_n.
\]

A linear dependence relation has the form
\vspace{-0.1cm}
\[
c_1\vect{a}_1+
c_2\vect{a}_2+
\cdots+
c_n\vect{a}_n
=
\mathbf{0},
\]

where at least one coefficient is non-zero.

But this can be written as
\vspace{-0.1cm}
\[
A
\begin{bmatrix}
c_1\\
c_2\\
\vdots\\
c_n
\end{bmatrix}
=
\mathbf{0}.
\]

Therefore,
\vspace{-0.1cm}
\[
\boxed{
\text{Non-trivial null vectors}
\iff
\text{Linear dependence among columns}.
}
\]

\end{frame}


% =========================================================
% TRIVIAL NULL SPACE
% =========================================================

\begin{frame}{Linear Independence and the Null Space}

The columns of $A$ are linearly independent if

\[
A\vect{x}=\mathbf{0}
\]

has only the trivial solution

\[
\boxed{
\vect{x}=\mathbf{0}.
}
\]

Thus,

\[
\boxed{
\operatorname{Null}(A)=\{\mathbf{0}\}
\iff
\text{columns of }A\text{ are linearly independent}.
}
\]

\vspace{0.4cm}

A non-zero vector in the null space therefore provides
a direct certificate of linear dependence.

\end{frame}


% =========================================================
% COLUMN SPACE AND NULL SPACE
% =========================================================

\begin{frame}{Column Space vs. Null Space}

For a matrix $A$:

\vspace{0.3cm}

\begin{columns}[T,totalwidth=\textwidth]

\begin{column}{0.44\textwidth}

\begin{block}{Column Space}

\[
\operatorname{Col}(A)
\]

Contains all possible outputs

\[
A\vect{x}.
\]

Its dimension is the rank.

\end{block}

\end{column}

\begin{column}{0.44\textwidth}

\begin{block}{Null Space}

\[
\operatorname{Null}(A)
\]

Contains all inputs satisfying

\[
A\vect{x}=\mathbf{0}.
\]

Its dimension is the nullity.

\end{block}

\end{column}

\end{columns}

\vspace{0.4cm}

One describes possible \textbf{outputs}; the other describes
inputs that are \textbf{collapsed to zero}.

\end{frame}


% =========================================================
% FUNDAMENTAL CONNECTION
% =========================================================

\begin{frame}{The Fundamental Connection}

For
\vspace{-0.1cm}
\[
A\vect{x}=\vect{b},
\]

two fundamental questions are:
\vspace{-0.1cm}
\[
\boxed{
\vect{b}\in\operatorname{Col}(A)\;?
}
\]

determines whether a solution exists.

If a solution $\vect{x}_p$ exists, then every solution has the form
\vspace{-0.1cm}
\[
\boxed{
\vect{x}
=
\vect{x}_p+\vect{x}_n,
\qquad
\vect{x}_n\in\operatorname{Null}(A).
}
\]

Therefore:
\vspace{-0.1cm}
\[
\boxed{
\text{Particular solution}
+
\text{Null-space solution}
}
\]

describes the complete solution set.

\end{frame}


% =========================================================
% INVERTIBLE MATRIX CONNECTION
% =========================================================

\begin{frame}{Column Space, Null Space and Invertibility}

For an $n\times n$ matrix $A$, the following are equivalent:

\[
\boxed{
\begin{array}{c}
A\text{ is invertible}\\
\Updownarrow\\
\operatorname{rank}(A)=n\\
\Updownarrow\\
\operatorname{Null}(A)=\{\mathbf{0}\}\\
\Updownarrow\\
\text{columns of }A\text{ are linearly independent}\\
\Updownarrow\\
\det(A)\neq0
\end{array}
}
\]

These equivalences connect the ideas developed throughout
linear algebra.

\end{frame}